% !PRINTED 413
\weberpage{413}
%<<CONT
$p^\pi$ ist, ein Theiler von $p^{\pi-1}$; d. h. jede nicht primitive Einheitswurzel vom Grade $p^\pi$ ist zugleich Einheitswurzel vom Grade $p^{\pi-1}$. Da es aber $p^\pi$ Einheitswurzeln vom Grade $p^\pi$ und nur $p^{\pi-1}$ Einheitswurzeln vom Grade $p^{\pi-1}$ giebt, so müssen
\[
p^\pi - p^{\pi-1} = p^\pi \left(1 - \frac{1}{p}\right)
\]
primitive Einheitswurzeln des Grades $p^\pi$ vorhanden sein. Daraus erhält man nach (5) die Anzahl aller primitiven $n^{\text{ten}}$ Einheitswurzeln
\weq{6}{\varphi(n) = n\,\Pi\left(1 - \frac{1}{p}\right),}
worin das Productzeichen $\Pi$ sich auf alle von einander verschiedenen in $n$ aufgehenden Primzahlen bezieht. Nur für den Fall $n = 1$ passt die Formel (6) nicht mehr; in diesem Falle ist $\varphi(1) = 1$ zu setzen. Die Zahl $\varphi(n)$ ist also niemals gleich Null. Da es hiernach für jeden Grad $n$ wenigstens eine primitive Einheitswurzel $r$ giebt, so lassen sich alle $n^{\text{ten}}$ Einheitswurzeln durch die Potenzen von $r$ darstellen:
\weq{7}{1, r, r^2, r^3 \ldots r^{n-1}.}
Ist $r^k$ irgend eine Potenz von $r$, so wird dann und nur dann $r^{km} = 1$ sein, wenn $km$ durch $n$ theilbar ist. Ist also $n = n'n''$ und $n'$ der grösste gemeinsame Theiler von $k$ und $n$, so muss $m$ durch $n''$ theilbar sein und $n''$ ist der Exponent der niedrigsten Potenz von $r^k$, die gleich 1 wird, d. h. $r^k$ ist eine primitive $n''^{\text{te}}$ Einheitswurzel. Es folgt hieraus der Satz:

\gesperrt{III. Ist $r$ primitive $n^{\text{te}}$ Einheitswurzel, so ist $r^k$ dann und nur dann primitive $n^{\text{te}}$ Einheitswurzel, wenn $k$ relativ prim zu $n$ ist.}

Nehmen wir $k$ aus der Reihe der Zahlen $1, 2, 3 \ldots n$, so \gesperrt{ergiebt sich der Satz der Zahlentheorie, dass $\varphi(n)$ gleich der Anzahl der Zahlen ist, die nicht grösser als $n$ und relativ prim zu $n$ sind.} Dies ist die ursprüngliche Definition des in der Zahlentheorie allgemein gebrauchten Zeichens $\varphi(n)$.

Ist $k$ relativ prim zu $n$, so kann man eine ganze Zahl $x$ so bestimmen, dass $kx \equiv 1\ (\mathrm{mod}\ n)$ wird. Sind dann $r, r_1$ zwei $n^{\text{te}}$ Einheitswurzeln, so kann nur dann $r^k = r_1^k$ sein, wenn $r = r_1$ ist, wie sich durch Erheben zur Potenz $x$ ergiebt. Daraus folgt nach III.:
%CONT>>
